\chapter{Discussion}

\section{Obtaining D-gene using the platform}
\subsection{Causative gene candidates list from variant analysis}
\hspace{1em}The variant layer showed a modest but statistically significant rank enrichment of the dPCD module (Mann--Whitney $p = 0.015$), with two cutoff-based over-representation checks pointing in the same direction without reaching significance on their own. This result, however, describes only an aggregate, module-wide trend: the 65 testable dPCD executor orthologs collectively skew toward higher damage scores than the rest of the scored gene set, which is a weaker and different claim than saying that any specific module gene -- let alone the true causative D-gene itself -- was individually prioritised.

\hspace{1em}Checking the true D-gene directly makes this distinction concrete. Fujimoto et al.\ \cite{fujimoto_transcriptional_2018} identify the sorghum D-gene as Sobic.006G147400 (LOC8084661 in the NCBI annotation used throughout this thesis), the ortholog of \textit{Arabidopsis} \textit{ANAC074}. This gene is nominally present in the candidate variant pool -- 575 variant--gene rows survive the SBC11-exclusion filter for it -- but its single most severe annotated consequence across the four accessions is a MODIFIER-impact \texttt{upstream\_gene\_variant}, the lowest severity band in the scoring scheme. Because gene-level scoring takes the maximum (worst) score across a gene's variants, LOC8084661 is pinned near the bottom of the ranked gene list despite technically clearing the initial site-selection filter. A structural-variant check at the same locus found no clean juicy-private candidate either: SVs overlapping LOC8084661 appear in several accession-sharing groups, including all four accessions together, but not in the SBC4\_SBC10\_SBC23 (juicy-only) group that would mirror the SNP/indel selection logic. Across every genomic track available in this four-accession panel, no variant offers a clean, phenotype-discriminating signal at the true D-gene locus.

\hspace{1em}This does not contradict the module-level enrichment result -- it illustrates that enrichment and single-gene recovery are separate claims that this analysis can satisfy independently of one another. A damaging-variant-based search will only surface a gene if its own coding sequence carries a qualifying lesion in this specific panel; if the causal difference instead lies in a regulatory element, a structural rearrangement not well captured by SNV/indel calling, or simply an allele not present among these four accessions (Fujimoto et al.\ report six independent naturally occurring nonfunctional D alleles across a broader sorghum diversity panel), a coding-impact score will not detect it.

\subsection{Causative gene candidates list from DMR analysis}
\hspace{1em}The DMR layer showed no enrichment of the dPCD module (hypergeometric ORA, fold $= 0.92$, $p = 0.836$), and the true D-gene offers a clean explanation for this null result rather than an unexplained failure. \texttt{LOC8084661} has zero DMR-feature rows anywhere in the annotated DMR table: DSS never called a differentially methylated region overlapping this locus in any of the six pairwise comparisons, so the gene was never even eligible to enter the DMR candidate pool, independent of the hypomethylation-in-SBC11 selection criterion applied afterwards.

\hspace{1em}This is consistent with the D-gene's known molecular basis. Fujimoto et al.\ report that the juicy phenotype in natural sorghum populations arises from independent nonfunctional \textit{D} alleles -- a genetic loss-of-function mechanism, not the epigenetic silencing of an otherwise intact wild-type allele. Promoter-methylation-based candidate selection is designed to detect the latter: a gene whose regulatory region has been differentially silenced between accessions that otherwise retain a functional coding sequence. Since the D locus's causal variation is allelic rather than epigenetic, a DMR-based search was targeting a molecular mechanism that does not apply at this specific locus, and the resulting null result is better read as a mismatch between the platform's epigenetic-silencing hypothesis and the D-gene's actual biology than as a failure of the DMR pipeline itself.

\subsection{D-gene transcription factor from the co-expression analysis}
\hspace{1em}The co-expression layer ranked the five NAC-family transcription factors present in the 66-gene dPCD module by how consistently each was co-expressed with the module's other 61 genes, nominating LOC8080837 (\textit{protein NTM1-like 9}, chromosome 7) as the most probable D-gene candidate. This is not the true D-gene: LOC8084661 sits on chromosome 6, matching Fujimoto et al.'s fine-mapping interval, and -- more fundamentally -- it is not a member of the 66-gene dPCD module at all. That module was built from orthologs of the conserved dPCD \textit{executor} gene families characterised by Olvera-Carrillo et al.\ \cite{olvera-carrillo_conserved_2015} (ribonucleases, serine carboxypeptidases, aspartic proteinases), whereas D acts upstream of these executors as their transcriptional switch, placing it outside the executor orthology group the module was curated from.

\hspace{1em}The co-expression search was therefore structurally unable to nominate the true D-gene, regardless of how well the ranking method performed among the five candidates it did consider: a guilt-by-association search that ranks a module's own members can only succeed if the true regulator is already a latent member of that module, and here it was excluded by the module's construction before any co-expression computation took place. This traces the platform's stated limitation for this case study -- failing to recover the D-gene via co-expression -- to a specific, identifiable cause: the module curation step, not the co-expression data or the ranking methodology, was the bottleneck.

\section{Obtaining aconitate isomerase using the platform}
\subsection{Causative gene candidates list}
\hspace{1em}Unlike the D-gene case study, aconitate isomerase has no independent ground truth to check candidate recovery against: the enzyme (EC 5.3.3.7) remains molecularly unidentified in any plant species, so there is no known gene analogous to LOC8084661 to search for among the variant-derived (21,928 scored genes) or DMR-derived (2,617 candidate genes) lists. In this setting the platform's output is necessarily an unvalidated hypothesis-generation exercise rather than a check against a known answer, and its credibility depends on internal consistency -- for instance, whether independent layers converge on the same genes -- rather than on recovering an established target.

\hspace{1em}Checking for such convergence, none of the co-expression layer's top five candidates (Table~\ref{tab:taa-coexpression-ranking}) appear among the DMR layer's 2{,}617 promoter-hypomethylated candidates. The absence of overlap is not, on its own, strong evidence against any individual layer: with candidate lists this large relative to the roughly 35,000-gene sorghum genome, substantial overlap would require real biological convergence to be likely by chance, and its absence here only weakens the additional claim that the layers agree with one another, not the validity of each layer in isolation.

\subsection{Identification of aconitate isomerase from the co-expression analysis}
\hspace{1em}The three-gene phytochrome-upregulated proxy module used to rank isomerase-family candidates was itself flagged, at the point it was constructed, as a lower-confidence substitute for a true literature-curated module: it reflects convergence on a qualitatively similar light-induction pattern across heterologous \textit{Arabidopsis} experiments rather than co-expression measured within a single shared dataset, unlike the dPCD module used for the D-gene case study. Any ranking built on it inherits that same uncertainty.

\hspace{1em}Consistent with this, none of the top five mean-MR-ranked candidates -- two inositol-3-phosphate synthase paralogs, a katanin p60 ATPase subunit, a DEAD-box RNA helicase, and achilleol B synthase -- carry a functional annotation connected to the TCA cycle or organic-acid metabolism, the pathway aconitate isomerase belongs to. This does not rule any of them out: enzyme-activity annotation in a non-model crop genome is itself incomplete, so a true aconitate isomerase could plausibly be mis-annotated under a different enzyme family. But it also means these candidates do not, on their own, tell an obvious mechanistic story the way a HIGH-impact coding variant overlapping a known dPCD executor gene would. They are best read as a prioritised hypothesis space for targeted follow-up (e.g.\ heterologous expression or knockout validation of the top-ranked genes) rather than as a confident nomination.

\hspace{1em}Taken together, the two case studies show that the platform can connect multi-omics evidence to candidate gene lists, and that each layer can produce a real, checkable signal when its underlying assumption matches the phenotype's true molecular mechanism. The D-gene case study also shows the reverse: a layer can fail cleanly and informatively when that assumption does not hold, as when a genetic loss-of-function mechanism is searched for with an epigenetic-silencing hypothesis, or when a regulatory gene is searched for within a module curated around its downstream executors. Applied to TAA, where no positive control exists to calibrate these assumptions against, the resulting candidate lists should be treated as the platform is designed to produce them: prioritised, testable hypotheses rather than final identifications.
